\Needspace{10\baselineskip}
\pdfbookmark[0]{Anzeigen. Mémorial du dépôt général de la guerre. 1830.}{anzeigen-memorial-1830}
\begin{center}
\rule{.85\textwidth}{.3pt}\par
Göttingische gelehrte Anzeigen.\quad 29. April 1830.\par
\rule{.85\textwidth}{.3pt}
\end{center}
\hspace*{\parindent}\textit{Mémorial du dépôt général de la guerre, imprimé par ordre du ministre.} Bd.\,I, 1829 (für 1802---1803), 696 Seiten. Bd.\,III, 1826 (für 1825), 466 Seiten. Bd.\,IV, 1828 (für 1826), 494 Seiten. Bd.\,V, 1829 (für 1827---1828), 490 Seiten im Quartformat. Mit zahlreichen Karten und Plänen. Paris bei \textsc{Picquet}.

Die Militärzeitschrift mit dem oben genannten Titel wurde 1802 auf Anregung des Generals \textsc{Andreossy}, des damaligen Direktors des Depots, gegründet. Man wollte nach und nach einen Teil der Schätze dieser bedeutenden Einrichtung, die während der Revolutionskriege ins Unermessliche bereichert worden war, für die Kriegswissenschaft und ihre Hilfswissenschaften, Geschichte, Topografie, Geodäsie, Statistik und so weiter, nutzbar machen. Dazu wollte man die Muße des gerade eingetretenen Friedens nutzen. So erschienen die ersten Hefte der Zeitschrift rasch nacheinander. Als der Krieg jedoch bald wieder ausbrach, geriet ihr Erscheinen allmählich erneut ins Stocken, und mit dem siebten
